
\clearpage
\section{Réalisations élémentaires du Modèle Standard}
\begingroup
\footnotesize
\setlength{\tabcolsep}{1.9pt}
\begin{longtable}{L{2.03cm}L{2.03cm}L{1.87cm}L{2.34cm}L{6.55cm}}
\toprule
\rowcolor{HMTNavy}\HMTtablehead{Particule} & \HMTtablehead{Charge} & \HMTtablehead{Spin} & \HMTtablehead{Statistique} & \HMTtablehead{Réalisation HMT--MD} \\
\midrule
\endfirsthead
\multicolumn{5}{l}{\sffamily\scriptsize\color{HMTGray}Suite du tableau}\\[-1pt]
\toprule
\rowcolor{HMTNavy}\HMTtablehead{Partícula} & \HMTtablehead{Carga} & \HMTtablehead{Espín} & \HMTtablehead{Estadística} & \HMTtablehead{Realización HMT--MD} \\
\midrule
\endhead
\midrule
\endfoot
\bottomrule
\endlastfoot
\(e^-\) & \(-1\) & \(1/\allowbreak{}2\) & fermiónica & punto fijo \((5,5)\), vacancia orientada \\
\rowcolor{HMTLight}\(\mu^-\) & \(-1\) & \(1/\allowbreak{}2\) & fermiónica & bloque leptónico de rango 8 \\
\(\tau^-\) & \(-1\) & \(1/\allowbreak{}2\) & fermiónica & bloque leptónico de rango 16 \\
\rowcolor{HMTLight}\(\nu_e\) & \(0\) & \(1/\allowbreak{}2\) & fermiónica & vector de interacción en soporte neutro 20 \\
\(\nu_\mu\) & \(0\) & \(1/\allowbreak{}2\) & fermiónica & vector de interacción en soporte neutro 20 \\
\rowcolor{HMTLight}\(\nu_\tau\) & \(0\) & \(1/\allowbreak{}2\) & fermiónica & vector de interacción en soporte neutro 20 \\
\(u\) & \(+2/\allowbreak{}3\) & \(1/\allowbreak{}2\) & fermiónica & rama up, generación 1, triplete de color \\
\rowcolor{HMTLight}\(d\) & \(-1/\allowbreak{}3\) & \(1/\allowbreak{}2\) & fermiónica & rama down, generación 1, triplete de color \\
\(c\) & \(+2/\allowbreak{}3\) & \(1/\allowbreak{}2\) & fermiónica & rama up, generación 2, triplete de color \\
\rowcolor{HMTLight}\(s\) & \(-1/\allowbreak{}3\) & \(1/\allowbreak{}2\) & fermiónica & rama down, generación 2, triplete de color \\
\(t\) & \(+2/\allowbreak{}3\) & \(1/\allowbreak{}2\) & fermiónica & rama up, generación 3, triplete de color \\
\rowcolor{HMTLight}\(b\) & \(-1/\allowbreak{}3\) & \(1/\allowbreak{}2\) & fermiónica & rama down, generación 3, triplete de color \\
\(\gamma\) & \(0\) & \(1\) & bosónica & sección nula, dos polarizaciones \\
\rowcolor{HMTLight}\(g^a\) & \(0\) & \(1\) & bosónica & adjunta de color, ocho generadores \\
\(W^\pm\) & \(\pm1\) & \(1\) & bosónica & bloque vectorial cargado conjugado \\
\rowcolor{HMTLight}\(Z^0\) & \(0\) & \(1\) & bosónica & bloque vectorial neutral \\
\(H\) & \(0\) & \(0\) & bosónica & sección escalar persistente \\
\end{longtable}
\endgroup
Las antipartículas se obtienen por \(C_M\). Los antiquarks pertenecen al triplete conjugado; \(W^-\) es el conjugado de \(W^+\); fotón, \(Z\), Higgs y gluones poseen las conjugaciones propias de sus representaciones.

El soporte B1 exacto de las clases elementales tiene rangos

\[
 1,\quad8,\quad16,\quad
 3\times20,\quad
 3\times16,\quad
 3\times20.
\]
Le nombre d'incidences classe--parcours est donc

\[
 1+8+16+60+48+60=193.
\]
Les incidences ne sont pas 193 parcours distincts : les trois saveurs neutriniques partagent le support neutre de rang vingt ; les saveurs up réutilisent un support de rang seize et les saveurs down, un support de rang vingt. Le projecteur de saveur doit raffiner ces supports. Dans l'union des quarks à 36 cellules, la couleur est exactement

\[
 12_R+12_G+12_B.
\]
Cette égalité certifie le support chromatique global, mais n'autorise pas encore à choisir une cellule par saveur.

\Needspace{7\baselineskip}
\section{Décomposition du secteur leptonique par génération, charge, spin et caractère de masse}
\Needspace{7\baselineskip}
\subsection{Électron : centre, lacunes et orientation}
L'électron est exceptionnel parce que sa fibre contient l'unique point fixe \((5,5)\). Sa réduction à une masse scalaire ne nécessite pas de briser une orbite. Le parcours actif est

\[
 e_{\rm act}=(2,2,5,-4,-3,-2)^2,
\]
tandis que l'enveloppe historique

\[
 e_{\rm env}=(4,3,5,-4,-3,-5)^2
\]
conserve le signe et le maximum absolu des deux feuillets. L'enveloppe retient le mot

\[
 \tau_e=+++---+++---,
\]
mais ne remplace pas l'histoire active.

Quatre données doivent demeurer distinctes :

\[
 \text{signature brute}=(24,0,12,0,-12),
\]
\[
 \text{signature paire CT--108}=(-12,-4,12,-6,12),
\]
\[
 \text{vecteur d'événements}=e_{\rm act},
\qquad
 \text{origine affine}=v_e=0.
\]
Le sélecteur de base \((-22,+15)\) possède deux lectures internes complémentaires. La décomposition régionale donne

\[
 -22=-(27-5),\qquad +15=5\cdot3=\binom62,
\]
et le cadre des lacunes

\[
 V_e=
 \begin{pmatrix}
 21&22\\
 6&7
 \end{pmatrix}
\]
satisfait

\[
 \det V_e=21\cdot7-22\cdot6=15.
\]
Le \(22\) mesure le déficit orienté de l'extension ; le \(15\), la rétention d'aire qui survit à la clôture. La particule électronique apparaît ainsi comme persistance centrale d'une lacune orientée, et non comme une masse insérée dans la cellule centrale.

L'équation complète est

\[
 \mathfrak e_e^\beta=
 \frac{\sqrt3}{4}
 \left[
 \exp\!\left(\frac{\varphi}{\pi^2}\right)
 -22\alpha_{\rm HMT}^{3}
 \right]
 (1+15\Delta_4)
 R_{\rm act}^{80-54\alpha_{\rm HMT}+6\Delta_4}.
\]
Dans la carte électronique,

\[
 \mathfrak e_e^\beta
 =0.5109989506900008086082571648\ldots\ {\rm MeV}.
\]
\Needspace{7\baselineskip}
\subsection{Réalisation du muon comme section leptonique de 2e génération : fibre, signature et caractère}
Le muon réside dans une multisection leptonique de rang huit. Son empreinte structurelle retrouvée comprend les cellules

\[
 \{21,23,25,39,43,57,59,61\},
\]
les dénombrements de rotation/changement \(22/25\), l'axe nord--sud,

\[
 (N,E,S,O)=(62,25,7,14),
\qquad
 (h_{369},r_5)=(57,9),
\]
et la coordonnée centrale

\[
 (n_A,s_C)=(42,-3).
\]
L'objet propre est l'opérateur comprimé

\[
 \widehat M_\mu=P_\mu\widehat M P_\mu.
\]
Sa réduction à un nombre exige de prouver que le projecteur de saveur et de génération commute avec l'opérateur, ou que l'état préparé a une variance nulle. Étant donnée la signature enrichie

\[
 (42,-3,8,0,2),
\]
l'évaluation du caractère est

\[
 105.658360294435829\ldots\ {\rm MeV}.
\]
Cette valeur est un contrôle exact de l'évaluation de cette signature ; la dérivation physique complète doit reproduire la signature à partir de l'enregistrement généalogique sélectionné par le couplage. La largeur muonique s'obtient à partir de l'opérateur temporel de désintégration, et non de \(\widehat M_\mu\).

\Needspace{7\baselineskip}
\subsection{Réalisation du tau comme section leptonique de 3e génération : fibre, signature et caractère}
Le tau occupe une multisection leptonique de rang seize. Son empreinte comprend

\[
 17/24,\qquad \text{axe est--ouest},\qquad
 (N,E,S,O)=(29,55,7,17),
\]
\[
 (h_{369},r_5)=(59,6),\qquad
 (n_A,s_C)=(64,0).
\]
La signature enrichie

\[
 (64,0,25,-1,6)
\]
produit

\[
 1776.8609176314323\ldots\ {\rm MeV}.
\]
Comme pour le muon, la structure établie est le bloc opératoriel ; la masse scalaire s'obtient lorsque le couplage tau sélectionne un bloc scalaire ou un pôle isolé.

\Needspace{7\baselineskip}
\subsection{Fibre neutrinique commune et séparation des saveurs par le lecteur PMNS}
Les trois saveurs neutriniques partagent l'espace

\[
 \mathcal H_\nu=
 \bigoplus_{j=1}^{4}\mathcal H_{\nu,G_j},
\qquad
 \dim\mathcal H_\nu=2+4+6+8=20.
\]
Les vecteurs

\[
 f_e,\quad f_\mu,\quad f_\tau
\]
constituent une base d'interaction ; ce ne sont pas trois vecteurs propres de masse. La profondeur \(G_j\) n'est pas non plus une quatrième saveur stérile : une interprétation stérile exige un projecteur physique supplémentaire.

L'extracteur neutre se définit par

\[
 T_\nu(m)=\sum_{t\in E_m}
 \omega(t)\varepsilon(t)\theta(t),
\qquad
 e_m^\nu=T_\nu(m)-T_\nu(m+6),
\]
\[
 \tau_m^\nu=\operatorname{sgn}(e_m^\nu).
\]
On en obtient

\[
 b_{120}^\nu
 =\sum_{a=0}^{3}
 \operatorname{sgn}\!\left(\sum_{m\in I_a}e_m^\nu\right),
\]
\[
 b_\zeta^\nu=\sum_m\tau_m^\nu\tau_{m+3}^\nu.
\]
La signature neutrinique prend la forme

\[
 v_\nu^{(i)}
 =(0,0,k_{\Delta_4}^{(i)},b_{120}^{\nu_i},b_\zeta^{\nu_i}).
\]
L'inversion bidirectionnelle change \(b_{120}\) et conserve \(b_\zeta\) ; cette asymétrie produit intérieurement les deux orientations de hiérarchie.

\Needspace{7\baselineskip}
\subsection{Nature fermionique du neutrino}
La neutralité ne détermine pas la statistique. Le secteur neutrinique se situe dans la complétion extérieure et porte l'holonomie spinorielle :

\[
 U_\nu(2\pi)=-I,\qquad U_\nu(4\pi)=I,
\]
\[
 \{a_{\nu,i},a_{\nu,j}^\dagger\}=\delta_{ij}.
\]
Chaque mode neutrinique physique est donc fermionique. La question Dirac/Majorana est ultérieure : elle détermine si la conjugaison échange deux sections distinctes ou identifie une section à sa conjuguée au sein d'une représentation réelle. Aucune des deux possibilités ne transforme le neutrino en boson.

\Needspace{7\baselineskip}
\subsection{Opérateur de masse et PMNS}
Soit \(P_\nu\) le projecteur de rang trois sur le sous-espace des modes massiques. On diagonalise d'abord

\[
 P_\nu\widehat M_\nu P_\nu
 =\sum_{i=1}^{3}m_i\,|n_i\rangle\langle n_i|.
\]
On construit ensuite le transport entre bases :

\[
 (f_e,f_\mu,f_\tau)
 =(n_1,n_2,n_3)U_{\rm PMNS}^\dagger.
\]
Un noyau de rang plein lisant feuillet, orientation, retenue, lacune et holonomie peut s'écrire

\[
 G_{ab}=
 \frac1{\sqrt{|L_a||N_b|}}
 \sum_{c\in L_a}\sum_{d\in N_b}K_{\rm reg}(c,d),
\]
\[
 U_{\rm PMNS}=\operatorname{polar}(G),
\qquad \det G\ne0.
\]
La phase interne de rang un produit les angles

\[
 (19.6807^\circ,56.4970^\circ,29.6224^\circ)
\]
et ne peut donc constituer le noyau PMNS physique complet. Le rang plein est une exigence mathématique, et non une retouche numérique.

Les observables

\[
 m_1,m_2,m_3,\qquad
 \Delta m_{ij}^2,\qquad
 m_\beta,\qquad
 m_{\beta\beta},\qquad
 \sum_i m_i
\]
sont des applications différentes du même spectre. Elles ne doivent pas être confondues dans une unique « masse du neutrino ».

\Needspace{7\baselineskip}
\section{Quarks, couleur, générations et mélange CKM}
\Needspace{7\baselineskip}
\subsection{Position structurelle des \(6\) saveurs de quarks}
Chaque quark se définit par

\[
 \mathfrak q_X=
 (\text{branche de charge},\text{génération},
 \mathbb Z_3^{\rm color},\text{multisection},
 \text{registre généalogique},\text{feuillet}).
\]
La position structurelle des six saveurs est résumée par leurs empreintes antérieures à toute carte de masse :

\Needspace{7\baselineskip}
\begingroup
\fontsize{7.2}{8.2}\selectfont
\setlength{\tabcolsep}{1.9pt}
\begin{longtable}{L{0.94cm}L{2.81cm}L{2.03cm}L{1.87cm}L{3.28cm}L{2.18cm}L{1.87cm}}
\toprule
\rowcolor{HMTNavy}\HMTtablehead{Saveur} & \HMTtablehead{Branche/\allowbreak{}génération} & \HMTtablehead{Giro/\allowbreak{}cambio} & \HMTtablehead{Eje} & \HMTtablehead{\((N,E,S,O)\)} & \HMTtablehead{\((h_{369},r_5)\)} & \HMTtablehead{\((n_A,s_C)\)} \\
\midrule
\endfirsthead
\multicolumn{7}{l}{\sffamily\scriptsize\color{HMTGray}Suite du tableau}\\[-1pt]
\toprule
\rowcolor{HMTNavy}\HMTtablehead{Sabor} & \HMTtablehead{Rama/\allowbreak{}generación} & \HMTtablehead{Giro/\allowbreak{}cambio} & \HMTtablehead{Eje} & \HMTtablehead{\((N,E,S,O)\)} & \HMTtablehead{\((h_{369},r_5)\)} & \HMTtablehead{\((n_A,s_C)\)} \\
\midrule
\endhead
\midrule
\endfoot
\bottomrule
\endlastfoot
\(u\) & positiva/\allowbreak{}1 & \(21/\allowbreak{}30\) & N--S & \((34,31,31,12)\) & \((51,13)\) & \((11,7)\) \\
\rowcolor{HMTLight}\(d\) & negativa/\allowbreak{}1 & \(20/\allowbreak{}32\) & E--O & \((27,35,18,28)\) & \((47,17)\) & \((17,8)\) \\
\(c\) & positiva/\allowbreak{}2 & \(21/\allowbreak{}29\) & E--O & \((20,38,24,26)\) & \((49,16)\) & \((61,8)\) \\
\rowcolor{HMTLight}\(s\) & negativa/\allowbreak{}2 & \(17/\allowbreak{}27\) & N--S & \((50,32,12,14)\) & \((57,11)\) & \((41,-3)\) \\
\(t\) & positiva/\allowbreak{}3 & \(21/\allowbreak{}26\) & N--S & \((46,27,16,19)\) & \((57,9)\) & \((100,-1)\) \\
\rowcolor{HMTLight}\(b\) & negativa/\allowbreak{}3 & \(21/\allowbreak{}30\) & E--O & \((12,63,9,24)\) & \((47,7)\) & \((71,-5)\) \\
\end{longtable}
\endgroup
Estas posiciones no son seis números elegidos: son seis clases de representación que actúan sobre tripletes de color. La ruta física puede ser una órbita o un bloque de rutas; no tiene que ser un punto.

\Needspace{7\baselineskip}
\subsection{Invarianza escalar del carácter de masa bajo la acción de color}
Sea \(P_q\) el proyector de sabor y generación y \(P_{\rm RGB}\) el proyector a la órbita de color. La masa quark es

\[
 \widehat M_q=
 P_qP_{\rm RGB}\widehat M P_{\rm RGB}P_q.
\]
Si

\[
 [\widehat M_q,U(g)]=0,\qquad g\in SU(3),
\]
alors

\[
 \widehat M_q=m_qI_3
\]
sur le triplet irréductible. Cette égalité explique pourquoi la saveur porte une masse commune bien que ses trois couleurs occupent des positions distinctes dans l'atlas.

\Needspace{7\baselineskip}
\subsection{Carte du schéma et de l'échelle}
La coordonnée HMT interne doit être transportée vers une définition renormalisée :

\[
 m_q^{\mathcal S}(\mu)
 =Z_m^{\mathcal S}(\mu,\mu_0)\,
 \mathcal R_q(\widehat M_q),
\]
avec la loi de composition

\[
 Z_m(\mu_2,\mu_0)
 =Z_m(\mu_2,\mu_1)Z_m(\mu_1,\mu_0).
\]
Pour \(u,d,s\), la comparaison naturelle utilise \(\overline{\rm MS}\) à \(2\,{\rm GeV}\) ; pour \(c,b\), l'échelle de la masse renormalisée elle-même ; pour \(t\), il faut déclarer si la carte est de pôle, MSR ou cinématique. Sans ce transport, une différence numérique n'est pas encore un résidu de la loi.

\Needspace{7\baselineskip}
\subsection{Évaluations des signatures déclarées}
Les coordonnées centrales historiques produisent :

\Needspace{7\baselineskip}
\begingroup
\footnotesize
\setlength{\tabcolsep}{2.1pt}
\begin{longtable}{L{4.68cm}L{10.61cm}}
\toprule
\rowcolor{HMTNavy}\HMTtablehead{Saveur} & \HMTtablehead{Évaluation} \\
\midrule
\endfirsthead
\multicolumn{2}{l}{\sffamily\scriptsize\color{HMTGray}Suite du tableau}\\[-1pt]
\toprule
\rowcolor{HMTNavy}\HMTtablehead{Sabor} & \HMTtablehead{Evaluación} \\
\midrule
\endhead
\midrule
\endfoot
\bottomrule
\endlastfoot
\(u\) & \(2.165924\allowbreak{}536173\allowbreak{}907\ldots\ {\rm MeV}\) \\
\rowcolor{HMTLight}\(d\) & \(4.679550\allowbreak{}162948\allowbreak{}874\ldots\ {\rm MeV}\) \\
\(c\) & \(1.270500\allowbreak{}838641\allowbreak{}911\ldots\ {\rm GeV}\) \\
\rowcolor{HMTLight}\(s\) & \(92.940093\allowbreak{}907845\allowbreak{}64\ldots\ {\rm MeV}\) \\
\(t\) & \(172.597479\allowbreak{}544019\allowbreak{}1\ldots\ {\rm GeV}\) \\
\rowcolor{HMTLight}\(b\) & \(4.190119\allowbreak{}763299\allowbreak{}790\ldots\ {\rm GeV}\) \\
\end{longtable}
\endgroup
Estas evaluaciones comprueban el carácter una vez dada la firma. La deducción completa exige obtener cada firma, su refinamiento de torre y su carta de esquema a partir del proyector de sabor y del registro genealógico correspondiente.

\Needspace{7\baselineskip}
\subsection{Determinación de los ángulos CKM desde la reflexión racional de las rutas quark}
Sean

\[
 A=1000\alpha_{\rm HMT},\qquad
 h=\frac{C^*}{6},\qquad
 \Delta=\frac{180}{\pi}\Delta_4.
\]
La grammaire angulaire donne

\[
 \begin{pmatrix}
 \theta_{12}\\
 \theta_{23}\\
 \theta_{13}
 \end{pmatrix}
 =
 \begin{pmatrix}
 2&-5&28\\
 0&7&-25/2\\
 1&-20&-2/3
 \end{pmatrix}
 \begin{pmatrix}
 A\\h\\\Delta
 \end{pmatrix},
\qquad
 \delta_{\rm CP}=9A.
\]
Le critère linéaire exact de réfutation est

\[
 3857\,\delta_{\rm CP}
 =9\bigl(
 1528\theta_{12}
 +3380\theta_{23}
 +801\theta_{13}
 \bigr).
\]
À partir des trois angles et de \(\delta_{\rm CP}\), on construit la matrice unitaire \(V_{\rm CKM}\). Sa fonction dans la loi des masses n'est pas de produire des valeurs propres, mais de transporter les opérateurs down vers la base up :

\[
 B_d^{(u)}=V_{\rm CKM}B_dV_{\rm CKM}^*.
\]
\Needspace{7\baselineskip}
\subsection{Invariant de Jarlskog}
L'invariant de violation CP est

\[
 J_{\rm CKM}
 =c_{12}c_{23}c_{13}^2
 s_{12}s_{23}s_{13}\sin\delta_{\rm CP},
\]
et l'évaluation interne est

\[
 J_{\rm CKM}
 =3.1189723326632947\times10^{-5}.
\]
De manière équivalente, la compatibilité entre spectres et mélange peut se lire dans le déterminant du commutateur des opérateurs hermitiens up/down. Une valeur non nulle de \(J_{\rm CKM}\) enregistre une orientation physique non éliminable ; ce n'est pas un correcteur de masse ajouté à chaque quark.

\Needspace{7\baselineskip}
\section{Secteurs de jauge et scalaire}
\Needspace{7\baselineskip}
\subsection{Carte du secteur nul : parcours de masse \(0\), représentations de jauge et persistance}
Soit

\[
 \mathcal A_{\rm sp}=\mathbb R[R]/(R^2-1),
 \qquad
 P_\pm=\frac{1\pm R}{2}.
\]
Pour

\[
 x=r_+P_++r_-P_-,
\]
la norme déployée est

\[
 N_{\rm sp}(x)=r_+r_-.
\]
La section nulle est une direction de norme nulle. Sa masse est exactement

\[
 m=0,
\]
sans passer par une limite de caractères positifs et sans attribuer \(\kappa=0\). La nullité appartient au type de la section.

\Needspace{7\baselineskip}
\subsection{Réalisation du photon dans le secteur nul électromagnétique : hélicité, jauge et caractère}
Le photon est la réalisation vectorielle transverse d'une section de jauge nulle. Le projecteur physique doit laisser exactement deux polarisations :

\[
 \operatorname{rg}P_\gamma^{\rm trans}=2,
\qquad
 P_\gamma^{\rm trans}k=0.
\]
La masse au repos est

\[
 m_\gamma=0.
\]
La propagation bidirectionnelle peut porter une anisotropie orientée

\[
 \epsilon=q_-^{90}-q_+^{90},
\qquad
 c_\pm=\frac{c}{1\pm\epsilon},
\]
tandis que la moyenne harmonique conserve

\[
 \frac{2}{1/c_++1/c_-}=c.
\]
Ainsi, l'aller et le retour peuvent différer sans modifier l'invariant two--way. Ce biais n'est pas une masse de photon.

Une réalisation Proca/Stückelberg distincte apparaît si un gap spectral de l'enregistrement généalogique satisfait

\[
 \omega^2=k^2+m_T^2
\]
et s'il existe un entrelaceur dimensionnel

\[
 m_T\longmapsto m_\gamma.
\]
Le retard induit obéirait à

\[
 \frac{\Delta T}{T}
 \simeq\frac12
 \left(\frac{m_\gamma c^2}{\hbar\omega}\right)^2.
\]
Le secteur Maxwell nul et la réalisation Proca massive sont deux branches distinctes ; l'une n'invalide pas l'autre.

\Needspace{7\baselineskip}
\subsection{Réalisation chromodynamique du secteur nul : octet de gluons et représentation adjointe de \(SU(3)\)}
Le trit résiduel engendre l'espace de couleur, dont la partie sans trace réalise

\[
 \mathfrak{sl}_3\longrightarrow\mathfrak{su}_3.
\]
Le gluon occupe la représentation adjointe de dimension huit. Tant que la symétrie de couleur demeure exacte,

\[
 m_g=0.
\]
Confinement, écrantage et gap spectral collectif sont des propriétés du spectre interactif ; ils ne constituent pas une masse de pôle élémentaire du gluon. Une masse effective doit préciser l'observable, la jauge, le régime et l'opérateur qui la définissent.

\Needspace{7\baselineskip}
\subsection{Réalisation des bosons \(W^\pm\) : charge, polarisation et caractère électrofaible}
Le secteur \(W\) est un bloc vectoriel chargé avec conjugaison

\[
 C_M:W^+\longleftrightarrow W^-.
\]
Son empreinte historique est

\[
 15/37,\qquad \text{axe est--ouest},\qquad
 (n_A,s_C)=(94,-1),
\]
et la signature enrichie déclarée est

\[
 (94,-1,0,-2,4).
\]
L'évaluation correspondante est

\[
 80.37896490970455\ldots\ {\rm GeV}.
\]
La réalisation complète est le pôle de la résolvante du bloc \(P_W\widehat M P_W\) ; sa partie imaginaire détermine la largeur. La dégénérescence de \(W^+\) et de \(W^-\) procède de la conjugaison de l'opérateur.

\Needspace{7\baselineskip}
\subsection{Réalisation du boson \(Z^0\) : mélange neutre, polarisation et caractère électrofaible}
Le \(Z^0\) est le bloc vectoriel neutre du couplage électrofaible. Son empreinte historique est

\[
 20/23,\qquad \text{axe est--ouest},\qquad
 (n_A,s_C)=(95,-1),
\]
et la signature enrichie

\[
 (95,-1,-11,0,-4).
\]
L'évaluation du caractère donne

\[
 91.18753344431341\ldots\ {\rm GeV}.
\]
Le projecteur neutre doit être construit avant la carte de pôle ; masse et largeur sont deux coordonnées du même pôle, et non une seule valeur numérique.

\Needspace{7\baselineskip}
\subsection{Réalisation du secteur de Higgs : mode scalaire, couplage et caractère de masse}
Le Higgs est une section scalaire persistante du bloc de jauge, de parité CP déclarée. Son empreinte est

\[
 24/29,\qquad \text{axe nord--sud},\qquad
 (n_A,s_C)=(97,9).
\]
L'évaluation conditionnelle du caractère produit

\[
 125.2924660425361\ldots\ {\rm GeV}.
\]
La définition HMT n'identifie pas « scalaire » à « masse » : le spin zéro fixe la représentation des rotations ; la masse procède du spectre du bloc scalaire couplé. La largeur et les canaux de désintégration exigent l'opérateur temporel.

\Needspace{7\baselineskip}
\subsection{Le moment magnétique anomal du muon}
La fonctionnelle interne

\[
 a_\mu=
 \frac{\alpha_{\rm HMT}}{2\pi}
 +\frac{\alpha_{\rm HMT}(\varphi-1)}{1000}
 +\frac{\alpha_{\rm HMT}^2}{1000(54+1)}
\]
donne

\[
 a_\mu=0.001165920713008014\ldots,
\]
et

\[
 g_\mu=2(1+a_\mu)
 =2.002331841426016028\ldots.
\]
Les trois termes correspondent respectivement à la rotation de base, au retour orienté dans la complétion \(1000\) et à la correction quadratique médiée par le défaut \(54\). Cet observable vérifie la structure interne du secteur muonique, mais ne sélectionne pas, à lui seul, son parcours de masse.

\Needspace{7\baselineskip}
\section{Construction des particules composées par composition non additive des parcours et des liaisons}
\Needspace{7\baselineskip}
\subsection{Opérateur de composition}
La composition des parcours n'est pas strictement additive, car frontière, ordre temporel et retenue interagissent. Pour deux enregistrements,

\[
 \mathcal L_1\star\mathcal L_2
 =\mathcal L_1+\mathcal L_2
 +\mathfrak b(\mathcal L_1,\mathcal L_2),
\]
où \(\mathfrak b\) est le cocycle de liaison. L'associativité exige

\[
 \mathfrak b(L_1,L_2)
 +\mathfrak b(L_1\star L_2,L_3)
 =
 \mathfrak b(L_2,L_3)
 +\mathfrak b(L_1,L_2\star L_3).
\]
Le lecteur de masse agit après cette composition :

\[
 (\Gamma_i,\mathcal T,\partial)
 \xrightarrow{\mathfrak C_{12}}
 \mathcal L_{\rm comp}
 \xrightarrow{L_{\rm tick}}
 \sigma_{\rm comp}
 \xrightarrow{\chi_{\rm full}}
 \widehat M_{\rm comp}.
\]
\Needspace{7\baselineskip}
\subsection{Réalisation mésonique par composition quark–antiquark et loi cocyclique de liaison}
Le domaine mésonique contient 432 expressions orientées

\[
 M(q,q')=q\otimes(q')^\vee
\]
à couleur compatible. Le projecteur physique impose le singulet de couleur et fixe

\[
 (J^{PC},I,I_3,S,C,B,\ldots),
\]
avec l'arbre temporel de liaison. La masse d'un méson n'est pas \(m_q+m_{\bar q}\) ; c'est une valeur propre de l'enregistrement composite.

Les six évaluations historiques de contrôle comprennent

\[
 m_{\pi^0}=134.9767243278721\ldots\ {\rm MeV},
\]
\[
 m_{\pi^\pm}=139.5704851715722\ldots\ {\rm MeV},
\]
\[
 m_{K^\pm}=493.6770856495306\ldots\ {\rm MeV},
\]
\[
 m_{K^0}=497.6111864977505\ldots\ {\rm MeV}.
\]
Ces valeurs sont des évaluations exactes d'enregistrements déclarés. La clôture généalogique état par état exige de régénérer chaque enregistrement au moyen de \(\mathfrak C_{12}\), y compris sa frontière de liaison.

\Needspace{7\baselineskip}
\subsection{Réalisation baryonique par composition trilinéaire des parcours de quarks et neutralisation de couleur}
Le domaine baryonique contient 1728 expressions RGB. Le projecteur physique doit imposer simultanément :

\begin{enumerate}
\item le singulet de couleur ;
\item la statistique fermionique ;
\item la saveur et l'isospin ;
\item le spin et la symétrisation ;
\item l'arbre triple avec retenue et mémoire.
\end{enumerate}
Pour le proton et le neutron, les signatures historiques

\[
 \sigma_p=(59,0,9,1,0),\qquad
 \sigma_n=(59,1,-34,0,1)
\]
donnent

\[
 m_p=938.2717515469849\ldots\ {\rm MeV},
\]
\[
 m_n=939.5658104725070\ldots\ {\rm MeV}.
\]
La précision apparente ne remplace pas la dérivation de l'enregistrement généalogique de liaison. Le signe et la grandeur de

\[
 m_n-m_p
\]
doivent découler de l'orientation, de l'isospin, des lacunes et de la frontière de liaison avant l'ouverture de la comparaison externe.

\Needspace{7\baselineskip}
\subsection{États excités et résonances hadroniques}
Une identité composée peut posséder plusieurs pôles. L'état physique est individualisé par

\[
 (\text{constituants},\mathcal T,\partial,
 J^{PC},I,\text{radial},\text{orbital},\text{canal}).
\]
La même composition de saveurs ne suffit pas à distinguer les excitations. Chaque pôle

\[
 z_k=M_k-\frac{i}{2}\Gamma_k
\]
doit procéder d'un vecteur propre de l'opérateur composé et d'un canal de désintégration. Cela étend la loi de quelques masses stables à tout le spectre des mésons et des baryons, sans transformer le catalogue externe en générateur.

\Needspace{7\baselineskip}
\subsection{Inventaire exhaustif des largeurs et des taux de transition des réalisations de l'atlas}
Soit \(P\) le sous-espace préparé et \(Q=I-P\). Le Hamiltonien effectif est

\[
 H_{\rm eff}(z)
 =PHP+PHQ(z-QHQ)^{-1}QHP.
\]
Les pôles de

\[
 \det(z-H_{\rm eff}(z))=0
\]
déterminent masses et largeurs. Le caractère de masse fournit la partie réelle de base de l'opérateur ; le terme d'autoénergie de \(Q\) fournit déplacement et désintégration. Les 384 largeurs de l'inventaire exigent donc un lecteur de canaux supplémentaire au lecteur des 471 masses.

\Needspace{7\baselineskip}
\section{Secteurs topologiques et particules hypothétiques}
\Needspace{7\baselineskip}
\subsection{Secteur topologique de Fibonacci : anneau basé, fusion et représentation de tresses}
L'incidence surfacique--volumétrique

\[
 F_{\rm av}=
 \begin{pmatrix}
 0&1\\1&1
 \end{pmatrix}
\]
réalise l'anneau

\[
 \tau^2=1+\tau,\qquad
 \operatorname{FPdim}\tau=\varphi.
\]
Le résultat primaire est une règle de fusion et une dimension quantique. Une particule anyonique de Fibonacci exige en outre des matrices \(F\) et \(R\), un pentagone, un hexagone, une chiralité, un Hamiltonien à gap spectral et une application de l'état TPK vers l'espace physique. Ce n'est qu'alors que l'on peut définir

\[
 m_{\rm brecha espectral}=\frac{\Delta E}{c_{\rm eff}^2}.
\]
\Needspace{7\baselineskip}
\subsection{Secteur d'Ising–Majorana : objets simples, règles de fusion et représentation de tresses}
La catégorie d'Ising satisfait

\[
 \psi^2=1,\qquad
 \psi\sigma=\sigma,\qquad
 \sigma^2=1+\psi,\qquad
 d_\sigma=\sqrt2.
\]
Quatre niveaux doivent être distingués :

\begin{enumerate}
\item l'involution des parcours \(C_M=-\operatorname{rev}\) ;
\item la parité de la complétion CAR ;
\item un spineur relativiste autoconjugué ;
\item un mode zéro topologique d'Ising.
\end{enumerate}
Un fermion relativiste de Majorana exige une représentation spinorielle réelle, une dynamique et un terme de masse. Un mode zéro de Majorana se caractérise par un gap spectral, une séparation de niveaux et une protection ; il ne possède pas de masse universelle déductible du seul anneau de fusion. L'identité d'un parcours avec son conjugué est une condition nécessaire, mais non suffisante, de la réalisation de Majorana.

\Needspace{7\baselineskip}
\subsection{Représentation de tresses sur \(\mathbb Z/6\mathbb Z\) et caractères topologiques}
Soit

\[
 w_6:B_n\longrightarrow\mathbb Z/6,
\qquad
 \chi_k(b)=e^{2\pi ikw(b)/6}.
\]
Les représentations

\[
 \rho_q(\sigma_i)=qP_i
\]
descendent au groupe symétrique seulement lorsque \(q^2=1\). Les quatre autres valeurs sont véritablement tressées. Leur existence mathématique n'implique pas quatre particules : une réalisation physique exige un Hamiltonien, un gap spectral, des opérateurs de tresse, une stabilité et un entrelaceur avec une multisection.

\Needspace{7\baselineskip}
\subsection{Impossibilité d'un tachyon persistant}
La droite d'action est positive,

\[
 \hbar_{\rm int}>0,
\]
le caractère satisfait

\[
 \chi_{\rm full}>0,
\]
et l'opérateur bêta s'obtient par conjugaison positive :

\[
 \widehat M^\beta
 =R_{\rm act}^{\widehat K/2}
 \widehat M^{(0)}
 R_{\rm act}^{\widehat K/2}\ge0.
\]
Par conséquent,

\[
 \operatorname{Spec}(\widehat M^\beta)\subset[0,\infty).
\]
Une section persistante stable ne peut engendrer

\[
 m^2<0.
\]
Si la linéarisation d'un état produit une courbure négative ou une fréquence imaginaire, cet état est une instabilité qui quitte la section, et non une particule supraluminique. Si le prolongement se termine, c'est une section virtuelle à obstruction terminale. Ainsi, « tachyon » se décompose en deux objets HMT correctement typés --- instabilité ou virtualité --- et ne constitue pas une espèce persistante.

\Needspace{7\baselineskip}
\subsection{Gravitation sans particule élémentaire obligatoire de spin \(2\)}
La réalisation gravitationnelle suit la chaîne

\[
 \text{holonomie de route}
 \longrightarrow(e,\omega)
 \longrightarrow(T,R)
 \longrightarrow
 S_{\rm PCH},
\]
où \(e\) est la costructure, \(\omega\) la connexion,

\[
 T=de+\omega\wedge e,\qquad
 R=d\omega+\omega\wedge\omega,
\]
et \(S_{\rm PCH}\) est l'action de Palatini--Cartan--Holst. La gravitation apparaît comme réponse globale de torsion et de courbure du réseau persistant. La vibration collective de l'univers modèle peut posséder une composante tensorielle de spin deux, mais HMT n'a pas besoin de postuler une particule élémentaire supplémentaire pour transporter la gravitation.

La constante gravitationnelle interne est typée par

\[
 G_{\rm HMT}
 =\frac{c_{\rm int}^3}{\hbar_{\rm int}}
 \bigl(\delta_\theta\Lambda_G\bigr)^2,
\]
où \(\delta_\theta\) est la torsion angulaire élémentaire et \(\Lambda_G\) la longueur spectrale sélectionnée par l'holonomie gravitationnelle. Cette équation convertit torsion et échelle de parcours en couplage gravitationnel sans utiliser de masse comme normalisation. La vibration totale est représentée par le Hessien

\[
 \mathcal H_{\rm univ}
 =\operatorname{Hess}_{(e,\omega)}S_{\rm PCH}
\]
sur l'état global. Ses modes sont des excitations collectives de l'univers modèle ; une composante tensorielle n'est pas automatiquement une particule séparée.

Une excitation appelée « graviton » doit démontrer trois propriétés :

\begin{enumerate}
\item un sous-espace tensoriel persistant ;
\item un pôle isolé de l'opérateur linéarisé ;
\item un couplage physique qui ne double pas la torsion déjà présente.
\end{enumerate}
Sans ces trois conditions, l'objet correct est le mode collectif de la géométrie, et non une nouvelle ligne de masse.

\Needspace{7\baselineskip}
\subsection{Barbero–Immirzi et la liaison hexade–octade}
Avec

\[
 A=1000\alpha_{\rm HMT},
\qquad
 q_\pm=\exp\!\left[-\frac{\pi}{180}(A\pm C^*)\right],
\]
la fonctionnelle hexade--octade utilise deux incidences distinctes des tours temporelles. La combinaison

\[
 \varepsilon_{\rm BI}
 =12\Delta S_{90}-\Delta S_{120}
\]
et le terme angulaire produisent

\[
 \gamma_{\rm BI}
 =\frac{\pi}{180}AC^*
 +\varepsilon_{\rm BI}
 =0.2740076726944593\ldots.
\]
Le poids \(12:-1\) procède de la lecture des douze secteurs dodécaphasiques
et de l'unique couture orientée de retour. L'équation
\(4a-3b=45\) vérifie ensuite que la paire \((12,1)\) est la solution
entière positive minimale. Cet objet ne doit pas être confondu avec l'extracteur dodécaphasique de masse ni avec le résidu angulaire \(90\)--\(120\) ; il partage l'opérateur harmonique, mais possède un autre domaine et une autre fonction.

\Needspace{7\baselineskip}
\subsection{Extensions de l'atlas vers les secteurs sombres : domaines, lecteurs et statut physique}
Un secteur sombre HMT est une multisection persistante avec

\[
 P_{\rm vis}P_{\rm dark}=0
\]
pour les couplages électromagnétique et chromatique, mais avec une incidence gravitationnelle ou de frontière possible :

\[
 P_{\rm grav}P_{\rm dark}\ne0.
\]
Son « obscurité » est une propriété du noyau de couplage, et non une étiquette de masse. La loi prédit d'abord un opérateur et son spectre ; seul un canal physique détermine si le mode est stable, métastable ou résonant. Cette définition permet de rechercher des multisections invisibles sans inventer une masse à partir de la valeur observée d'une densité cosmologique.

\Needspace{7\baselineskip}
\subsection{Section neutrinique stérile : compatibilité de fibre, mélange et absence de charge de jauge}
Un neutrino stérile n'est pas la profondeur \(G_4\). C'est un vecteur propre \(n_s\in\mathcal H_\nu\) qui satisfait

\[
 P_{\rm weak}n_s=0,\qquad
 P_{\rm mass}n_s=n_s.
\]
Il peut se mélanger aux trois modes actifs au moyen d'un noyau neutre étendu, mais son existence exige

\[
 \dim\ker(P_{\rm weak}|_{\mathcal H_\nu})>0
\]
et une valeur propre de masse dans ce noyau. Cette définition produit un critère interne de recherche : diagonaliser d'abord le bloc neutre, puis vérifier son couplage faible.

\Needspace{7\baselineskip}
\subsection{Sections axionique et pseudoscalaires : caractère topologique et couplage de bord}
Un candidat axionique est une section pseudoscalaire dont l'orientation change sous CP,

\[
 CP\,a\,CP^{-1}=-a,
\]
et dont le déplacement compense une phase topologique du secteur de couleur. Dans HMT, il doit correspondre à un mode impair de torsion/holonomie et à un projecteur qui entrelace ce mode avec la densité topologique de jauge. Sa masse s'obtient à partir de la dérivée seconde du potentiel effectif,

\[
 m_a^2=
 \left.\frac{\partial^2V_{\rm eff}}{\partial a^2}\right|_{a=a_0},
\]
une fois \(V_{\rm eff}\) construit à partir de l'enregistrement généalogique des lacunes et de la frontière. La seule existence d'un canal \(C^*\mapsto-C^*\) ne suffit pas à affirmer l'existence d'une particule axionique.

\Needspace{7\baselineskip}
\subsection{Photon sombre et vecteur massif caché}
Un photon sombre exige une seconde section de jauge \(A'\) et un couplage de mélange

\[
 \mathcal L_{\rm mix}
 =-\frac{\varepsilon_{\rm kin}}2F_{\mu\nu}F'^{\mu\nu}.
\]
En termes de fibres, \(\varepsilon_{\rm kin}\) est un entrelaceur entre deux sous-espaces nuls. Si le second secteur acquiert un gap spectral de Stückelberg ou de frontière,

\[
 \widehat M_{A'}^2\ne0,
\]
un vecteur massif sombre apparaît. HMT prédit que le paramètre de mélange doit procéder d'une incidence de mémoire commune ; il ne peut être fixé avec la masse que l'on cherche à expliquer.

\Needspace{7\baselineskip}
\subsection{WIMP et état sombre stable}
Un WIMP est une section sombre qui satisfait en outre :

\[
 [P_{\rm weak},P_X]\ne0,\qquad
 [P_{\rm em},P_X]=[P_{\rm color},P_X]=0,
\]
et possède un mode stable

\[
 |\lambda_X|=1.
\]
La stabilité peut être protégée par une parité de l'enregistrement généalogique. Sa masse est une valeur propre du bloc sombre faiblement couplé. « WIMP » ne constitue pas une nouvelle famille mathématique : c'est une condition conjointe de couplage et de persistance sur une multisection.

\Needspace{7\baselineskip}
\subsection{Section monopolaire : charge magnétique, condition de Dirac et compatibilité topologique}
Un monopôle se définit par une classe topologique non triviale du fibré de jauge :

\[
 \frac1{2\pi}\int_{S^2}F=n\ne0.
\]
HMT doit réaliser cet entier comme holonomie nette d'une famille fermée de parcours. Sa masse, si le mode est dynamique et stable, appartient au spectre de l'opérateur de défaut topologique. Une lacune ponctuelle ou une retenue non nulle ne constituent pas, à elles seules, un monopôle ; la descente à la classe entière doit être démontrée.

\Needspace{7\baselineskip}
\subsection{Extensions supersymétriques de l'atlas : appariement des représentations et des caractères}
Une supersymétrie HMT exigerait un opérateur impair \(Q\) échangeant les complétions extérieure et symétrique :

\[
 Q:\mathcal F_-\leftrightarrow\mathcal F_+,
\qquad
 Q^2=H,
\]
et entrelaçant la loi des masses,

\[
 Q\widehat M_-=\widehat M_+Q.
\]
Ce n'est qu'alors que les spectres bosonique et fermionique seraient appariés. L'existence séparée de CAR et de CCR ne démontre pas la supersymétrie et ne prédit pas, à elle seule, des sélectrons, neutralinos ou gluinos. Chaque partenaire exige un entrelaceur non nul et une règle de brisure déterminant sa séparation de masse.

\Needspace{7\baselineskip}
\subsection{Tétraquarks, pentaquarks, boules de glu et hybrides}
La grammaire de composition se prolonge sans changer de loi :

\[
 \mathfrak C_{12}^{(n)}:
 \Gamma_1\times\cdots\times\Gamma_n
 \times\mathcal T_n\times\partial_n
 \longrightarrow\mathcal L_{\rm comp}.
\]
Un tétraquark utilise \(qq\bar q\bar q\) ; un pentaquark, \(qqqq\bar q\) ; une boule de glu, des enregistrements de l'adjointe de couleur projetés sur le singulet ; un hybride, des parcours de quarks et de jauge dans un même arbre. La classification physique dépend du projecteur de couleur, du spin, de la parité et de l'arbre de liaison, et non de la seule liste des constituants.

\Needspace{7\baselineskip}
\subsection{Sections de dilaton et de radion : modes scalaires de volume et de compactification}
Le dilaton est un mode de la section d'échelle ; le radion, un mode qui modifie une longueur ou un volume interne. Dans la théorie des masses, ils apparaissent comme fluctuations de la carte et de la mémoire géométrique :

\[
 t_0\mapsto e^{\phi_d}t_0,\qquad
 \mathcal V_{\rm int}\mapsto e^{\phi_r}\mathcal V_{\rm int}.
\]
Pour devenir des particules, ils doivent posséder un terme cinétique, un potentiel, un projecteur persistant et un pôle. Tant qu'ils ne les possèdent pas, ce sont des degrés collectifs de la carte dimensionnelle, et non des espèces scalaires supplémentaires.

\Needspace{7\baselineskip}
\subsection{Exclusion des fantômes physiques}
La mesure spectrale d'un état physique est positive :

\[
 \mu_{P,a}^{\rm mass}(B)\ge0.
\]
Une direction de norme négative n'appartient pas à l'espace des états préparé. Si une paramétrisation introduit un « fantôme », la réalisation physique doit quotienter la redondance de jauge ou rejeter le couplage. Cette exclusion est indépendante de l'interdiction tachyonique : un fantôme viole la positivité de la norme ; un tachyon persistant violerait la positivité de l'opérateur de masse.

\Needspace{7\baselineskip}
\section{Inventaire universel des masses et des largeurs}
\Needspace{7\baselineskip}
\subsection{Correspondance ultérieure entre les réalisations HMT–MD et le dénombrement physique externe}
L'inventaire de comparaison contient

\[
 1170\ \text{identités},\qquad
 438\ \text{groupes d'identité},
\]
\[
 471\ \text{observables de masse},\qquad
 384\ \text{observables de largeur}.
\]
Sa partition est

\Needspace{7\baselineskip}
\begingroup
\footnotesize
\setlength{\tabcolsep}{2.1pt}
\begin{longtable}{L{5.04cm}L{5.04cm}L{5.04cm}}
\toprule
\rowcolor{HMTNavy}\HMTtablehead{Domaine} & \HMTtablehead{Masses} & \HMTtablehead{Largeurs} \\
\midrule
\endfirsthead
\multicolumn{3}{l}{\sffamily\scriptsize\color{HMTGray}Suite du tableau}\\[-1pt]
\toprule
\rowcolor{HMTNavy}\HMTtablehead{Dominio} & \HMTtablehead{Masas} & \HMTtablehead{Anchuras} \\
\midrule
\endhead
\midrule
\endfoot
\bottomrule
\endlastfoot
partículas generales y leptónicas & 53 & 9 \\
\rowcolor{HMTLight}quarks & 8 & 1 \\
mesones & 243 & 225 \\
\rowcolor{HMTLight}bariones & 166 & 149 \\
registro gravitatorio de contraste & 1 & 0 \\
\rowcolor{HMTLight}\textbf{Total} & \textbf{471} & \textbf{384} \\
\end{longtable}
\endgroup
Las 471 masas no son 471 especies elementales: incluyen estados excitados, resonancias, distintas cargas y múltiples realizaciones compuestas. Las 384 anchuras son observables temporales y no deben ser leídas por el carácter de masa.

\Needspace{7\baselineskip}
\subsection{Identidad interna \(\mathcal I_X\) de una realización material y criterio de individualización frente a observaciones múltiples}
Cada estado comparado debe poseer una identidad interna

\[
 \mathcal I_X=
 (\mathcal C_X,\mathcal R_X,\Gamma_X,\mathcal L_X,
 \sigma_X,\eta_X,P_X,\mathfrak M_X,\mathcal T_X),
\]
où :

\begin{itemize}
\item \(\mathcal C_X\) est le type physique de réalisation déterminé par l'objet, la représentation et le codomaine ;
\item \(\mathcal R_X\) est la représentation physique ;
\item \(\Gamma_X\) est un parcours, une orbite ou une composition orientée ;
\item \(\mathcal L_X\) est l'enregistrement généalogique ;
\item \(\sigma_X\) est la signature complète ;
\item \(\eta_X\) est la tour ;
\item \(P_X\) est le projecteur du couplage ;
\item \(\mathfrak M_X\) est le canal de mesure ;
\item \(\mathcal T_X\) est l'arbre de composition, le cas échéant.
\end{itemize}
Deux observations externes appartiennent au même état HMT lorsqu'elles partagent cette identité et ne diffèrent que par l'expérience, la carte ou l'estimateur. L'existence de deux mesures ne crée pas deux particules.

\Needspace{7\baselineskip}
\subsection{Morphisme de comparaison entre invariants internes et observables métrologiques}
La comparaison exige le diagramme

\[
 \mathcal L_M^{\rm HMT}
 \xrightarrow{\;\varsigma_{\mathcal S,\mu}\;}
 \mathbb R_{>0}u_M^{\mathcal S,\mu}
 \xleftarrow{\;\text{mesure}\;}
 \mathcal O_X^{\rm ext}.
\]
Une ligne scientifiquement complète contient :

\[
 (m_X^{\rm HMT},m_X^{\rm ext},
 u,\mathcal S,\mu,\Sigma_{\rm exp},\Sigma_{\rm HMT},
 \text{convention de pôle}).
\]
Pour les quarks, schéma et échelle sont obligatoires ; pour les résonances, la convention de pôle ; pour les limites, le niveau de confiance ; pour les observables corrélés, la matrice de covariance.

\Needspace{7\baselineskip}
\subsection{Évaluation explicite de \(17\) réalisations représentatives de l'atlas}
Les évaluations suivantes conservent la séparation entre la loi et la donnée de signature :

\Needspace{7\baselineskip}
\begingroup
\footnotesize
\setlength{\tabcolsep}{2.1pt}
\begin{longtable}{L{4.68cm}L{10.61cm}}
\toprule
\rowcolor{HMTNavy}\HMTtablehead{État} & \HMTtablehead{Valeur du caractère dans sa carte} \\
\midrule
\endfirsthead
\multicolumn{2}{l}{\sffamily\scriptsize\color{HMTGray}Suite du tableau}\\[-1pt]
\toprule
\rowcolor{HMTNavy}\HMTtablehead{État} & \HMTtablehead{Valeur du caractère dans sa carte} \\
\midrule
\endhead
\midrule
\endfoot
\bottomrule
\endlastfoot
\(\mu\) & \(105.658360\allowbreak{}294435\allowbreak{}829\ldots\ {\rm MeV}\) \\
\rowcolor{HMTLight}\(\tau\) & \(1776.860917\allowbreak{}631432\allowbreak{}3\ldots\ {\rm MeV}\) \\
\(u\) & \(2.165924\allowbreak{}536173\allowbreak{}907\ldots\ {\rm MeV}\) \\
\rowcolor{HMTLight}\(d\) & \(4.679550\allowbreak{}162948\allowbreak{}874\ldots\ {\rm MeV}\) \\
\(c\) & \(1.270500\allowbreak{}838641\allowbreak{}911\ldots\ {\rm GeV}\) \\
\rowcolor{HMTLight}\(s\) & \(92.940093\allowbreak{}907845\allowbreak{}64\ldots\ {\rm MeV}\) \\
\(t\) & \(172.597479\allowbreak{}544019\allowbreak{}1\ldots\ {\rm GeV}\) \\
\rowcolor{HMTLight}\(b\) & \(4.190119\allowbreak{}763299\allowbreak{}790\ldots\ {\rm GeV}\) \\
\(W\) & \(80.378964\allowbreak{}909704\allowbreak{}55\ldots\ {\rm GeV}\) \\
\rowcolor{HMTLight}\(Z\) & \(91.187533\allowbreak{}444313\allowbreak{}41\ldots\ {\rm GeV}\) \\
\(H\) & \(125.292466\allowbreak{}042536\allowbreak{}1\ldots\ {\rm GeV}\) \\
\rowcolor{HMTLight}\(\pi^0\) & \(134.976724\allowbreak{}327872\allowbreak{}1\ldots\ {\rm MeV}\) \\
\(\pi^\pm\) & \(139.570485\allowbreak{}171572\allowbreak{}2\ldots\ {\rm MeV}\) \\
\rowcolor{HMTLight}\(K^\pm\) & \(493.677085\allowbreak{}649530\allowbreak{}6\ldots\ {\rm MeV}\) \\
\(K^0\) & \(497.611186\allowbreak{}497750\allowbreak{}5\ldots\ {\rm MeV}\) \\
\rowcolor{HMTLight}\(p\) & \(938.271751\allowbreak{}546984\allowbreak{}9\ldots\ {\rm MeV}\) \\
\(n\) & \(939.565810\allowbreak{}472507\allowbreak{}0\ldots\ {\rm MeV}\) \\
\end{longtable}
\endgroup
L'électron bêta n'appartient pas à ce tableau, car son parcours central, ses coefficients de base et son lecteur complet sont déterminés par la construction elle-même.

\Needspace{7\baselineskip}
